% Part: first-order-logic
% Chapter: completeness
% Section: compactness-direct

\documentclass[../../../include/open-logic-section]{subfiles}

\begin{document}

\iftag{FOL}
      {\olfileid{fol}{com}{cpd}}
      {\olfileid{pl}{com}{cpd}}

\olsection{સઘનતા પ્રમેયની સીધી સાબિતી}

પૂર્ણતા પ્રમેયનો આશરો લીધા વિના, તેની સાબિતીમાંના વિચારો વાપરીને
સઘનતા પ્રમેય સીધો સાબિત કરી શકાય છે. પૂર્ણતા પ્રમેયની સાબિતીમાં આપણે
!!{sentence}sના સુસંગત ગણ~$\Gamma$થી શરૂ કરીને તેને !!{sentence}sના
સુસંગત\iftag{FOL}{, સંતૃપ્ત}{} અને !!{complete} ગણ~$\Gamma^*$ સુધી
વિસ્તાર્યો હતો. પછી~$\Gamma^*$માંથી રચેલા
\iftag{FOL}{પદ-નિદર્શ~$\Struct{M(\Gamma^*)}$}{!!{valuation}
~$\pAssign{v(\Gamma^*)}$}માં~$\Gamma$નાં બધાં !!{sentence}s સાચાં છે
તેમ બતાવ્યું હતું; એટલે~$\Gamma$ સંતોષ્ય હતો.

વાક્યોનો સાન્ત રીતે સંતોષ્ય ગણ સંતોષ્ય છે તેમ બતાવવા આ જ રીત વાપરી
શકાય છે. સત્યતા સહાયક પ્રમેય સુધી લઈ જતાં પરિણામોનાં એવાં અનુરૂપ
સ્વરૂપો જ સાબિત કરવાં પડશે જેમાં ``સુસંગત''ને બદલે ``સાન્ત રીતે
સંતોષ્ય'' મૂકીએ.

\begin{prop}
\ollabel{prop:fsat-ccs}
ધારો કે~$\Gamma$ !!{complete} અને સાન્ત રીતે સંતોષ્ય છે. ત્યારે:
\begin{enumerate}
\tagitem{prvAnd}{$(!A \land !B) \in \Gamma$ જો અને માત્ર જો
  $!A \in \Gamma$ અને $!B \in \Gamma$ બંને.}{}

\tagitem{prvOr}{$(!A \lor !B) \in \Gamma$ જો અને માત્ર જો કાં તો
  $!A \in \Gamma$ અથવા $!B \in \Gamma$.}{}

\tagitem{prvIf}{$(!A \lif !B) \in \Gamma$ જો અને માત્ર જો કાં તો
  $!A \notin \Gamma$ અથવા $!B \in \Gamma$.}{}
\end{enumerate}
\end{prop}

\tagprob{FOL}
\begin{prob}
~$\Proves$ વાપર્યા વિના \olref[fol][com][cpd]{prop:fsat-ccs}
સાબિત કરો.
\end{prob}
\tagendprob

\tagprob{notFOL}
\begin{prob}
~$\Proves$ વાપર્યા વિના \olref[pl][com][cpd]{prop:fsat-ccs}
સાબિત કરો.
\end{prob}
\tagendprob

\iftag{FOL}{%
\begin{lem}
\ollabel{lem:fsat-henkin} દરેક સાન્ત રીતે સંતોષ્ય ગણ~$\Gamma$ને
સંતૃપ્ત અને સાન્ત રીતે સંતોષ્ય ગણ~$\Gamma'$ સુધી વિસ્તારી શકાય છે.
\end{lem}
}{}

\tagprob{FOL}
\begin{prob}
\olref[fol][com][cpd]{lem:fsat-henkin} સાબિત કરો. (સંકેત: નિર્ણાયક
પગલું !!{derivation}s અથવા સુસંગતતાનો જરાય આશરો લીધા વિના એ બતાવવાનું
છે કે જો~$\Gamma_n$ સાન્ત રીતે સંતોષ્ય હોય, તો
$\Gamma_n\cup\{!D_n\}$ પણ સાન્ત રીતે સંતોષ્ય છે.)
\end{prob}
\tagendprob

\iftag{FOL}{%
\begin{prop}\ollabel{prop:fsat-instances}
ધારો કે~$\Gamma$ પૂર્ણ, સાન્ત રીતે સંતોષ્ય અને સંતૃપ્ત છે.
\begin{tagenumerate}{prvEx,prvAll}
\tagitem{prvEx}{$\lexists[x][!A(x)]\in\Gamma$ જો અને માત્ર જો ઓછામાં
  ઓછા એક બંધ પદ~$t$ માટે $!A(t)\in\Gamma$.}{}
\tagitem{prvAll}{$\lforall[x][!A(x)]\in\Gamma$ જો અને માત્ર જો દરેક
  બંધ પદ~$t$ માટે $!A(t)\in\Gamma$.}{}
\end{tagenumerate}
\end{prop}
}{}

\tagprob{FOL}
\begin{prob}
\olref[fol][com][cpd]{prop:fsat-instances} સાબિત કરો.
\end{prob}
\tagendprob

\begin{lem}
\ollabel{lem:fsat-lindenbaum} દરેક સાન્ત રીતે સંતોષ્ય ગણ~$\Gamma$ને
!!a{complete} અને સાન્ત રીતે સંતોષ્ય ગણ~$\Gamma^*$ સુધી વિસ્તારી શકાય
છે.
\end{lem}

\tagprob{FOL}
\begin{prob}
\olref[fol][com][cpd]{lem:fsat-lindenbaum} સાબિત કરો. (સંકેત: નિર્ણાયક
પગલું એ બતાવવાનું છે કે જો~$\Gamma_n$ સાન્ત રીતે સંતોષ્ય હોય, તો
$\Gamma_n\cup\{!A_n\}$ અથવા $\Gamma_n\cup\{\lnot !A_n\}$માંથી કોઈ એક
સાન્ત રીતે સંતોષ્ય છે.)
\end{prob}
\tagendprob

\tagprob{notFOL}
\begin{prob}
\olref[pl][com][cpd]{lem:fsat-lindenbaum} સાબિત કરો. (સંકેત: નિર્ણાયક
પગલું એ બતાવવાનું છે કે જો~$\Gamma_n$ સાન્ત રીતે સંતોષ્ય હોય, તો
$\Gamma_n\cup\{!A_n\}$ અથવા $\Gamma_n\cup\{\lnot !A_n\}$માંથી કોઈ એક
સાન્ત રીતે સંતોષ્ય છે.)
\end{prob}
\tagendprob

\begin{thm}[સઘનતા]
\ollabel{thm:compactness-direct} $\Gamma$ સંતોષ્ય છે જો અને માત્ર જો તે
સાન્ત રીતે સંતોષ્ય છે.
\end{thm}

\begin{proof}
જો~$\Gamma$ સંતોષ્ય હોય, તો એવી
\iftag{FOL}{!!a{structure}~$\Struct{M}$}{!!a{valuation}~$\pAssign{v}$}
છે કે દરેક $!A\in\Gamma$ માટે
$\iftag{FOL}{\Sat{M}}{\pSat{v}}{!A}$. અલબત્ત, આ
\iftag{FOL}{$\Struct{M}$}{$\pAssign{v}$}~$\Gamma$ના દરેક સાન્ત ઉપગણને
પણ સંતોષે છે, તેથી~$\Gamma$ સાન્ત રીતે સંતોષ્ય છે.

હવે ધારો કે~$\Gamma$ સાન્ત રીતે સંતોષ્ય છે.\iftag{FOL}{
  \olref{lem:fsat-henkin} મુજબ એવો સાન્ત રીતે સંતોષ્ય, સંતૃપ્ત ગણ
  $\Gamma'\supseteq\Gamma$ છે.}{} \olref{lem:fsat-lindenbaum} મુજબ
~$\iftag{FOL}{\Gamma'}{\Gamma}$ને !!a{complete} અને સાન્ત રીતે સંતોષ્ય
ગણ~$\Gamma^*$ સુધી વિસ્તારી શકાય છે\iftag{FOL}{, અને~$\Gamma^*$ હજી
સંતૃપ્ત છે}{}. \olref[mod]{defn:termmodel}માંની કલમો પ્રમાણે
\iftag{FOL}{પદ-નિદર્શ~$\Struct{M(\Gamma^*)}$}{!!{valuation}
~$\pAssign{v(\Gamma^*)}$} રચો.\iftag{FOL}{ નોંધો કે
\olref[mod]{prop:quant-termmodel} એ~$\Gamma^*$ સુસંગત છે (કે તે
!!{complete} અથવા સંતૃપ્ત છે) એ વાત પર આધાર રાખતો નહોતો; તે ફક્ત
$\Struct{M(\Gamma^*)}$ આવૃત છે એ હકીકત પર આધાર રાખતો હતો.}{}
સત્યતા સહાયક પ્રમેય (\olref[mod]{lem:truth})ની સાબિતી ત્યારે પણ ચાલે
છે જ્યારે \olref[ccs]{prop:ccs}\iftag{FOL}{ અને
\olref[hen]{prop:saturated-instances}ના સંદર્ભોને અનુક્રમે
\olref{prop:fsat-ccs} અને \olref{prop:fsat-instances}ના સંદર્ભોથી}{}
બદલીએ. જ્યાં મૂળ સાબિતી
સુસંગતતાનો સીધો ઉપયોગ કરે છે ત્યાં સાન્ત સંતોષ્યતા એ જ જરૂરી હકીકતો
આપે છે: $\lfalse\notin\Gamma^*$, $\ltrue\in\Gamma^*$, અને
$\lnot !B\in\Gamma^*$ જો અને માત્ર જો $!B\notin\Gamma^*$; કારણ કે
આમાંથી વિરુદ્ધ કિસ્સામાં એક અથવા બે સભ્યોનો અસંતોષ્ય સાન્ત ઉપગણ મળે.
આથી~$\eq$ ન ધરાવતા દરેક $!A\in\Gamma$ને
$\Struct{M(\Gamma^*)}$ સાચું બનાવે છે.

\iftag{FOL}{જો~$\Gamma$માં~$\eq$ આવે, તો
\olref[ide]{defn:term-model-factor}ની જેમ~$\Struct{M(\Gamma^*)}$નું
ભાગફલ લો. જરૂરી સંબંધ~$\approx$ સ્વવાચક, સમમિત, પરંપરિત અને
વિધેયો તથા વિધેયસંકેતો સાથે સુસંગત છે: દરેક નિષ્ફળતા પૂર્ણતાથી તેનો
નિષેધ~$\Gamma^*$માં મૂકે અને પછી સાન્ત રીતે અસંતોષ્ય ઉપગણ આપે છે.
તેથી ભાગફલ સુવ્યાખ્યાયિત છે અને \olref[ide]{lem:truth}ની અનુમાનાત્મક
સાબિતી સાન્ત-સંતોષ્ય આવૃત્તિમાં પણ ચાલે છે. પરિણામે ભાગફલ નિદર્શ
~$\Gamma$ના દરેક વાક્યને સાચું બનાવે છે.}{}
\end{proof}

\sourcecorrection{OLCO-018}{મૂળની \texttt{compactness-direct.tex},
પંક્તિઓ~139--142માં સત્યતા સહાયક પ્રમેયમાંના ફક્ત બે સંદર્ભો બદલવાનું
કહ્યું છે. પરંતુ મૂળ સાબિતીના અસત્યતા, સત્યતા અને નિષેધના કિસ્સા પણ
સુસંગતતાનો સીધો ઉપયોગ કરે છે. અહીં એ ત્રણ ઉપયોગોને સાન્ત સંતોષ્યતાથી
કેવી રીતે બદલવા તે સ્પષ્ટ કર્યું છે.}

\sourcecorrection{OLCO-019}{મૂળની \texttt{compactness-direct.tex},
પંક્તિઓ~133--142માં માત્ર સાદો પદ-નિદર્શ અને
\olref[mod]{lem:truth} વપરાય છે, જ્યારે તે સત્યતા સહાયક પ્રમેય સ્પષ્ટ
રીતે~$\eq$ ન ધરાવતા વાક્યો પૂરતું છે. સામાન્ય FOL નિવેદન માટે જરૂરી
ભાગફલ પદ-નિદર્શનો કિસ્સો અહીં ઉમેર્યો છે; તેનો આધાર
\olref[ide]{defn:term-model-factor} અને \olref[ide]{lem:truth} છે.}

\sourcecorrection{OLCO-022}{મૂળની \texttt{compactness-direct.tex},
પંક્તિઓ~140--142માં \texttt{\string\iftag}નો ફરજિયાત ત્રીજો દલીલ-ભાગ
ખૂટે છે. અહીં પ્રથમ-ક્રમ શાખા પછી ખાલી ``ના'' શાખા~\texttt{\{\}}
ઉમેરી છે.}

\tagprob{FOL}
\begin{prob}
\olref[fol][com][cpd]{thm:compactness-direct}ની સાબિતી માટે જરૂરી
સ્વરૂપમાં સત્યતા સહાયક પ્રમેય
(\olref[fol][com][mod]{lem:truth})ની સંપૂર્ણ સાબિતી લખો.
\end{prob}
\tagendprob

\tagprob{notFOL}
\begin{prob}
\olref[pl][com][cpd]{thm:compactness-direct}ની સાબિતી માટે જરૂરી
સ્વરૂપમાં સત્યતા સહાયક પ્રમેય
(\olref[pl][com][mod]{lem:truth})ની સંપૂર્ણ સાબિતી લખો.
\end{prob}
\tagendprob

\end{document}
